\section{First Properties of Schemes}

\begin{exercise}\label{ex3.1}
	Let $V_i=\operatorname{Spec}\!\left(B_i\right)$ be an affine open cover, and let $f^{-1}\!\left(V_i\right)$ be covered by $U_{ij}=\operatorname{Spec}\!\left(A_{ij}\right)$, where $A_{ij}$ is a finitely generated $B_i$-algebra. Let $V_i\cap V=\bigcup_kD_V\!\left(b_{ik}\right)$ where $b_{ik}\in B$, then $D_V\!\left(b_{ik}\right)\hookrightarrow V_i$ induces a map $B_i\rightarrow B_{b_{ik}}$. Also, $U_{ij}\cap f^{-1}\!\left(D_V\!\left(b_{ik}\right)\right)=D_{U_{ij}}\!\left(f_{U_{ij}}^\#\left(b_{ik}\right)\right)$. Thus, $f^{-1}(V)$ can be covered by $\operatorname{Spec}\!\left(\left(A_{ij}\right)_{f_{U_{ij}}^\#\!\left(b_{ik}\right)}\right)$, where $\left(A_{ij}\right)_{f_{U_{ij}}^\#\!\left(b_{ik}\right)}=A_{ij}\otimes_{B_i}B_{b_{ik}}$ is a finitely generated $B_{b_{ik}}$-algebra.
\end{exercise}

\begin{exercise}\label{ex3.2}
	Use the expression $U_{ij}\cap f^{-1}\!\left(D_V\!\left(b_{ik}\right)\right)$ in \Cref{ex3.1}. Since $V=\bigcup_{i,k}D_V\!\left(b_{ik}\right)$ is quasi-compact, so $i$, $k$ can be chosen to have finitely many. $j$ has finitely many by hypothesis.
\end{exercise}

\begin{exercise}\label{ex3.3}
	\begin{enumerate}[label=(\alph*),ref=\theexercise(\alph*)]\crefalias{enumi}{exercise}
		\addtocounter{enumi}{2}
		\item\label{ex3.3.3} Let $U_i=\operatorname{Spec}\!\left(A_i\right)$ be an affine open cover of $f^{-1}(V)$ such that $A_i$ is a finitely generated $B$-algebra. Let $U_i\cap U=\bigcup_kD_U\!\left(a_{ik}\right)$, where $a_{ik}\in A$. Since $U$ is affine, the union $U=\bigcup_{i,k}D_U\!\left(a_{ik}\right)$ can be finite. Since each $D_U\!\left(a_{ik}\right)$ is affine, there are finitely many $a_{ikl}^\prime\in A_i$ such that $D_U\!\left(a_{ik}\right)=\bigcup_lD_{U_i}\!\left(a_{ikl}^\prime\right)$.
		
		The disjoint union of all inclusions $D_{U_i}\!\left(a_{ikl}^\prime\right)\subseteq D_U\!\left(a_{ik}\right)$ induces an injective map
		\begin{equation*}
			\varphi:A_{a_{ik}}\rightarrow\prod_l\left(A_i\right)_{a_{ikl}^\prime}\qquad \alpha\mapsto\left(\alpha|_{D_{U_i}\!\left(a_{ikl}^\prime\right)}\right)_l.
		\end{equation*}		
		However, if $\left(\beta_l\right)\in\prod_l\left(A_i\right)_{a_{ikl}^\prime}$ is an image from $A_{a_{ik}}$, the sheaf property on $D_U\!\left(a_{ik}\right)$ implies $\beta_{l_1}|_{D_{U_i}\!\left(a_{ikl_1}^\prime a_{ikl_2}^\prime\right)}=\beta_{l_2}|_{D_{U_i}\!\left(a_{ikl_1}^\prime a_{ikl_2}^\prime\right)}$, so $A_{a_{ik}}$ is isomorphic to a quotient ring of $\prod_l\left(A_i\right)_{a_{ikl}^\prime}$. The ring $\prod_l\left(A_i\right)_{a_{ikl}^\prime}$ is a finitely generated $B$-algebra, so is $A_{a_{ik}}$. By a similar method, $A$ is a finitely generated $B$-algebra.
	\end{enumerate}
\end{exercise}

\begin{exercise}\label{ex3.4}
	Similar to \Cref{ex3.1}.
\end{exercise}

\begin{exercise}\label{ex3.5}
	\begin{enumerate}[label=(\alph*),ref=\theexercise(\alph*)]\crefalias{enumi}{exercise}
		\item \textcolor{cyan}{Atiyah-Macdonald Exercise 8.4}.
		\item Let $V_i=\operatorname{Spec}\!\left(B_i\right)$, $f^{-1}\!\left(V_i\right)=\operatorname{Spec}\!\left(A_i\right)$. $B_i\rightarrow A_i$ is an integral homomorphism. If $C\subseteq X$ is closed, by \textcolor{cyan}{Atiyah-Macdonald Exercise 5.1}, $f\!\left(C\cap f^{-1}\!\left(V_i\right)\right)=f(C)\cap V_i$ is closed in subspace $V_i$, so $f(C)$ is closed.
		\item $k[x]\rightarrow k\!\left[x,x^{-1}\right]\times k$ defined to be $k[x]\hookrightarrow k\!\left[x,x^{-1}\right]$ and $x\mapsto0$ in $k[x]\rightarrow k$.
	\end{enumerate}
\end{exercise}

\begin{exercise}\label{ex3.6}
	By \Cref{3.1} and \Cref{ex2.9}, $\xi$ is unique. Also, $\xi$ is contained in every open subset. For the $U=\operatorname{Spec}A$, $A=\mathcal{O}_X(U)$ is an integral domain, so $\mathcal{O}_{X,\xi}=\left(\mathcal{O}_X|_U\right)_{\xi}=A_{(0)}$ is a field.
\end{exercise}

\begin{exercise}\label{ex3.7}
	Let $V=\operatorname{Spec}B\subseteq Y$ and $U_i=\operatorname{Spec}\!\left(A_i\right)\subseteq f^{-1}(V)$, such that $A_i$ is a finitely generated $B$-algebra. Since $X$ is irreducible, $U_i$ is dense in $X$, so $f\!\left(U_i\right)$ is dense in $V$, and by \textcolor{cyan}{Atiyah-Macdonald Exercise 1.21(v)}, $B\rightarrow A_i$ is injective.
	
	Let $S=B-\{0\}$. Then $S^{-1}A_i$ is a finitely generated $k(B)$-algebra. By \textcolor{cyan}{Noether's normalization lemma}, $S^{-1}A_i$ is integral over $k(B)\!\left[t_1,\dots,t_n\right]$. By \textcolor{cyan}{Atiyah-Macdonald Exercise 5.10}, $\operatorname{Spec}\!\left(S^{-1}A_i\right)\rightarrow\operatorname{Spec}\!\left(k(B)\!\left[t_1,\dots,t_n\right]\right)$ is surjective. Since $f$ is generically finite, $\operatorname{Spec}\!\left(S^{-1}A_i\right)=\operatorname{Spec}\!\left(A_i\otimes_Bk(B)\right)$ is a finite set, so $n=0$. Therefore, $S^{-1}A_i$ is a finite dimensional $k(B)$-vector space, so $S^{-1}A_i=k(A)$, which is a finite field extension of $k(B)$.
	
	Let $x_{ij}$ generate $A_i$ over $B$. Then $x_{ij}$ satisfies an algebraic dependence equation whose coefficients are all in $B$. Let $b_{ij}$ be the leading coefficient and let $s_i$ be their multiplication with respect to $j$. Then each $x_{ij}\in A_i\!\left[s_i^{-1}\right]$ are integral, so $A_i\!\left[s_i^{-1}\right]$ is a finitely generated $B\!\left[s_i^{-1}\right]$-module.
	
	Since $X$ is irreducible, $W=\bigcap_iU_i$ is a non-empty open set. There exists $a_i\in A_i\!\left[s_i^{-1}\right]$ such that $D\!\left(a_i\right)\in W$. Since $A_i\!\left[s_i^{-1}\right]$ is finite over $B\!\left[s_i^{-1}\right]$, take $b_i$ to be the constant coefficient of the monic equation, so $b_i\in\left(a_i\right)$, and therefore $D\!\left(b_i\right)\subseteq D\left(a_i\right)\subseteq W$. Let $b$ be the multiplication of $b_i$ with respect to $i$, then $U=D(b)$ is the desired open set.
\end{exercise}

\begin{exercise}\label{ex3.8}
	Let $V=\operatorname{Spec}B$ and let $x\in U\cap V$. Then there are $f\in A$ and $g\in B$ such that $D_V(g)\subseteq D_U(f)\subseteq U\cap V$ with $x\notin D_U(f)$ and $x\notin D_V(g)$. The restriction $D_U(f)\subseteq V$ induces a homomorphism $B\rightarrow A_f$ in which $g\mapsto a/f^n$. By \Cref{ex2.16.1}, $D_U(f)\cap D_V(g)=D_{\operatorname{Spec}\left(A_f\right)}\!\left(a/f^n\right)=D_U(af)$, so $D_V(g)=D_U(af)$. Therefore, $U\cap V$ can be covered by $\left\{W_i\right\}$ such that $W_i=D_U\!\left(f_i\right)=D_V\!\left(g_i\right)$. Since \textcolor{cyan}{localization commutes with integral closure}, the isomorphisms $\varphi_i:D_{\tilde{U}}\!\left(f_i\right)\rightarrow D_{\tilde{V}}\!\left(g_i\right)$ are compatible on overlap. Then use \Cref{ex2.12}.
	
	Break to affine case, we need to prove that the integral closure of a finitely generated $k$-algebra $A$ is a finitely generated $A$-module. By \textcolor{cyan}{Noether's normalization lemma}, $A$ is finite over $k\!\left[y_1,\dots,y_r\right]$. Let $K=\operatorname{Frac}\!\left(k\!\left[y_1,\dots,y_r\right]\right)$ and $L=\operatorname{Frac}A$. Then $L/K$ is a finite field extension. Let $B$ be the integral closure of $k\!\left[y_1,\dots,y_r\right]$ in $L$. \textcolor{red}{Generalization of} \textcolor{cyan}{Atiyah-Macdonald Exercise 5.15}.
\end{exercise}

\begin{exercise}\label{ex3.9}
	Use \Cref{3.3}. $k(s)\otimes_kk(t)$ is a PID. The prime ideals are $(0)$ and ideals generate by the elements which is irreducible and cannot be written in the form $f(s)\otimes_kg(t)$.
\end{exercise}

\begin{exercise}\label{ex3.10}
	\begin{enumerate}[label=(\alph*),ref=\theexercise(\alph*)]\crefalias{enumi}{exercise}
		\item\label{ex3.10.1} \textcolor{cyan}{Atiyah-Macdonald Exercise 3.21(iv)} is the affine case. Let $\left\{U_i\right\}$ be an affine open cover of $X$, then each $\varphi_i:\left(U_i\right)_y\rightarrow U_i\cap f^{-1}(y)$ is a homeomorphism. By the universal property of fiber product, the morphisms in the diagram
		\begin{equation*}
			\begin{tikzcd}
				\left(U_i\cap U_j\right)_y\arrow[r]\arrow[d]\arrow[rd]&\left(U_i\right)_y\arrow[d]\\
				\left(U_j\right)_y\arrow[r]&X_y
			\end{tikzcd}
		\end{equation*}
		are all unique, so the diagram commutes, and therefore $\varphi_i$ and $\varphi_j$ are compatible on overlap.
		\item\label{ex3.10.2} $X_y=\operatorname{Spec}\!\left(k[t]/\!\left(t^2-a\right)\right)$. $X_\eta=\operatorname{Spec}\!\left(k(s)[t]/\!\left(t^2-s\right)\right)$. $k$ is algebraically closed, so $t^2-s$ is irreducible.
	\end{enumerate}
\end{exercise}

\begin{exercise}\label{ex3.11}
	\begin{enumerate}[label=(\alph*),ref=\theexercise(\alph*)]\crefalias{enumi}{exercise}
		\item\label{ex3.11.1} Use \textcolor{cyan}{Atiyah-Macdonald Exercise 2.2} for the affine case. When only $X$ is affine, use a method similar to \Cref{ex3.10.1} to glue morphisms. Let $g:X^\prime\rightarrow X$, $X_i$ an affine cover of $X$, $Y_i=f^{-1}\!\left(X_i\right)$, and $X^\prime_i=g^{-1}\!\left(X_i\right)$. From the commutative diagram of $Y\times_XX^\prime_i$, the image of $Y\times_XX^\prime_i\rightarrow Y$ is inside $Y_i$. By the uniqueness from \Cref{3.3}, we have $Y\times_XX^\prime_i\cong Y_i\times_XX^\prime_i\cong Y_i\times_{X_i}X^\prime_i$. Then glue $X^\prime_i$.
		\item\label{ex3.11.2} Let $D(f)$ be a basic open set of $X$ and let $y\in Y\cap D(f)$, then $f\notin\mathfrak{p}_{i(y)}$. The map $i_y^\#:\mathcal{O}_{X,i(y)}\rightarrow\mathcal{O}_{Y,y}$ is a surjective local homomorphism, which implies $\left(i^\#(f)\right)_y=i_y^\#\!\left(f_{i(y)}\right)\notin\mathfrak{m}_{Y,y}$. By \Cref{ex2.16.1}, $U\cap Y\cap D(f)$ is a basic open set of $Y$ for each open affine $U\subseteq Y$. Since a closed subset of a quasi-compact space is quasi-compact, the cover of $Y$ can use only finitely many $f_i$. The set $\left\{f_i\right\}$ can be expanded so it generates the unit ideal of $A$. Restrict each $f_i$ on $Y$ and use \Cref{ex2.17.2}.
		\item\label{ex3.11.3} Let $U_i=\operatorname{Spec}\!\left(A_i\right)$ be an affine cover of $X$. By \hyperref[ex3.11.2]{(b)}, $Y^\prime\cap U_i=V\!\left(\mathfrak{a}_i\right)$; by hypothesis, $Y\cap U_i=V\!\left(\!\sqrt{\mathfrak{a}_i}\right)$. There is a natural homomorphism $A/\mathfrak{a}_i\rightarrow\left.A\middle/\!\sqrt{\mathfrak{a}_i}\right.$, which induces a morphism $\varphi_i:Y\cap U_i\rightarrow Y^\prime\cap U_i$. Therefore, we have a morphism of sheaves $\mathcal{O}_{U_i}|_{Y^\prime\cap U_i}\rightarrow\mathcal{O}_{U_i}|_{Y\cap U_i}$, which is exactly $\mathcal{O}_X|_{Y^\prime\cap U_i}\rightarrow\mathcal{O}_X|_{Y\cap U_i}$. Restrict on $U_i\cap U_j$, so $\varphi_i$ and $\varphi_j$ are compatible on overlap.
		\item\label{ex3.11.4} Let $\mathfrak{a}_i=\ker\!\left(A_i\rightarrow\Gamma\!\left(f^{-1}\!\left(U_i\right),\mathcal{O}_Z\right)\right)$, and let $Y_i=\operatorname{Spec}\!\left(A_i/\mathfrak{a}_i\right)$. By \Cref{ex2.4}, the ring homomorphism $A_i/\mathfrak{a}_i\rightarrow\Gamma\!\left(f^{-1}\!\left(U_i\right),\mathcal{O}_Z\right)$ induces a morphism $\varphi_i:f^{-1}\!\left(U_i\right)\rightarrow Y_i$. Combine $\varphi_i$ with $Y_i\rightarrow U_i$ and use \Cref{ex2.4} again to prove the corresponding diagram commutes. Combine $\varphi_i$ with $Y_i\rightarrow Y$ and use a method similar to \hyperref[ex3.11.3]{(c)}, $\varphi_i$ can be glued.
	\end{enumerate}
\end{exercise}

\begin{exercise}\label{ex3.12}
	\begin{enumerate}[label=(\alph*),ref=\theexercise(\alph*)]\crefalias{enumi}{exercise}
		\item Similar to \textcolor{cyan}{Atiyah-Macdonald Exercise 1.21(iv)}, use isomorphism $S/\ker\varphi\rightarrow T$.
		\item Use \Cref{ex2.14.3}.
	\end{enumerate}
\end{exercise}

\begin{exercise}\label{ex3.13}
	\begin{enumerate}[label=(\alph*),ref=\theexercise(\alph*)]\crefalias{enumi}{exercise}
		\item 
		\item 
		\item 
		\item 
		\item 
		\item 
		\item 
	\end{enumerate}
\end{exercise}

\begin{exercise}\label{ex3.14}
	\textcolor{cyan}{Atiyah-Macdonald Exercise 5.24(ii) and 5.26}. A counter example is $\mathbb{Z}$.
\end{exercise}

\begin{exercise}\label{ex3.15}
	\begin{enumerate}[label=(\alph*),ref=\theexercise(\alph*)]\crefalias{enumi}{exercise}
		\item 
		\item 
		\item 
	\end{enumerate}
\end{exercise}

\begin{exercise}\label{ex3.16}
	Suppose there is a closed subset such that $\mathcal{P}$ doesn't hold. Let $\Sigma$ be the collection of closed subsets of $X$ such that $\mathcal{P}$ doesn't hold, then $\Sigma$ is not empty. Since $X$ is Noetherian, $\Sigma$ has a minimal element $Z$. However, $\mathcal{P}$ holds for each closed subset of $Z$, which contradicts to the hypothesis.
\end{exercise}

\begin{exercise}\label{ex3.17}
	\begin{enumerate}[label=(\alph*),ref=\theexercise(\alph*)]\crefalias{enumi}{exercise}
		\item \Cref{ex2.9}.
		\item If $C\subseteq X$ is minimal closed, the closure of any point in $C$ is $C$ itself.
		\item The closure of a point $x$ is irreducible with $x$ the generic point. Prove that $x\in\overline{\{y\}}$ and $y\in\overline{\{x\}}$ simultaneously is impossible.
		\addtocounter{enumi}{2}
		\item 
	\end{enumerate}
\end{exercise}

\begin{exercise}\label{ex3.18}
	\begin{enumerate}[label=(\alph*),ref=\theexercise(\alph*)]\crefalias{enumi}{exercise}
		\item 
		\item 
		\item 
		\item 
	\end{enumerate}
\end{exercise}

\begin{exercise}\label{ex3.19}
	\begin{enumerate}[label=(\alph*),ref=\theexercise(\alph*)]\crefalias{enumi}{exercise}
		\item 
		\item 
		\item 
	\end{enumerate}
\end{exercise}

\begin{exercise}\label{ex3.20}
	\begin{enumerate}[label=(\alph*),ref=\theexercise(\alph*)]\crefalias{enumi}{exercise}
		\item 
		\item 
		\item 
		\item 
		\item 
		\item 
	\end{enumerate}
\end{exercise}

\begin{exercise}\label{ex3.21}
	
\end{exercise}

\begin{exercise}\label{ex3.22}
	\begin{enumerate}[label=(\alph*),ref=\theexercise(\alph*)]\crefalias{enumi}{exercise}
		\item 
		\item 
		\item 
		\item 
		\item 
	\end{enumerate}
\end{exercise}

\begin{exercise}\label{ex3.23}
	
\end{exercise}
